\documentclass[10pt, a4paper, oneside, titlepage]{article}

\usepackage[utf8]{inputenc}
\usepackage[T1]{fontenc}
\usepackage{lmodern, textcomp}
\usepackage[english, french]{babel}
\usepackage{csquotes}

\usepackage[activate={true,nocompatibility}, final, babel=true,
% 	tracking=true, spacing=true, %kerning=true,
	factor=1100, stretch=10, shrink=10]{microtype}

\usepackage[heightrounded, top=4cm, bottom=3.5cm, left=3.5cm, right=3.5cm,
			marginparsep=0.5cm%, marginparwidth=4cm
			]{geometry}
% default top=4.5cm, bottom=4cm, left=4.5cm, right=4.5cm
\usepackage{changepage}
\usepackage{marginnote}

\usepackage{eurosym, xspace}
\usepackage[nottoc]{tocbibind}
\usepackage{graphicx, subfig, float}
\usepackage[multiple]{footmisc}
\usepackage{authblk}
\usepackage[backend=bibtex8, citestyle=numeric-comp, maxbibnames=99, sorting=nyt,
			sortcites=true, firstinits=true]{biblatex}
% TODO: améliorer le style
/home/harold/texmf/tex/latex/biblatex/escrime.def

\usepackage{amsmath, amsfonts, amssymb, siunitx}

% \usepackage{makeidx}
\usepackage[cc]{titlepic}
\usepackage{xr-hyper}

\usepackage[usenames, dvipsnames]{xcolor}
\usepackage{framed, mdframed}
\usepackage[amsmath, hyperref, framed]{ntheorem}
% \usepackage{amsthm, thmtools}


% hyperfootnotes=false
\usepackage[pdftex, colorlinks=true, breaklinks=true,
	urlcolor=blue, linkcolor=blue, citecolor=red,
	pdfcreator={PDFLaTeX}, pdfauthor={Harold Erbin}
]{hyperref}
\usepackage[all]{hypcap}

\usepackage[sort&compress, nameinlink, french]{cleveref}


\usepackage{escrime}


\hypersetup{
	pdftitle={Arts martiaux historiques - ateliers},
	% pdfkeywords={Mots clés PDF},
	% pdfsubject={Sujet PDF}
}


\pagestyle{plain}
\graphicspath{{images/}}


\renewcommand{\Affilfont}{\small}
\newcommand{\email}[1]{\thanks{\href{mailto:#1}{\nolinkurl{#1}}}}
% \sisetup{group-separator = {,}, separate-uncertainty=true}
% enable non-breakable space
\DeclareUnicodeCharacter{00A0}{~}


\addbibresource{art_martial}
\iffalse
\bibliography{art_martial}
\fi

\externaldocument{amhe_recueil}


\newcommand{\auteur}[1]{\noindent Auteur : #1 \bigskip}
\newcommand{\base}[1]{\noindent Basé sur : #1 \bigskip}


\title{Arts martiaux historiques -- ateliers}

\author[*]{Harold Erbin\email{harold.erbin@gmail.com}}
\affil[*]{Chapitre des armes, Paris, France}
\affil[*]{Club d'escrime ancienne, École Normale Supérieure, Paris, France}

\titlepic{\includegraphics[width=0.8\textwidth]{logo_cda.pdf}}


\begin{document}

\thispagestyle{empty}
\maketitle


\clearpage
\setcounter{page}{2}

\thispagestyle{empty}
\begin{center}
	\setlength{\baselineskip}{1.5\baselineskip}{\fontsize{12}{14}\selectfont{}
	Copyright : Les ateliers restent la propriété de leurs auteurs respectifs. \\
	Le reste du texte est diffusé sous la licence libre \\
	\textbf{Licence Art Libre} :} \\
	\url{http://artlibre.org/licence/lal/}
\end{center}

\vspace{\stretch{2}}

\noindent
% {\setlength{\tabcolsep}{0.5ex}
\begin{tabular}{ll}
	Sites du projet : &
		\url{https://github.com/melsophos/manuel-amhe} \\
		&
		\url{https://chapitredesarmes.wordpress.com/travaux} \\
	Téléchargement : & 
		\url{https://www.dropbox.com/s/1dpoy35uxz1gaj2/amhe_ateliers.pdf}
\end{tabular}
% }

% \vspace{\stretch{1}}


\clearpage
\pdfbookmark[1]{\contentsname}{toc}
\microtypesetup{protrusion=false}
\tableofcontents
\microtypesetup{protrusion=true}

% \begin{center}
	
% \end{center}

\newpage

\section*{Introduction}


Ce document contient diverses idées d'ateliers.
Chacun fait appel à un ensemble de notions et est composé d'exercices de préparation suivis des techniques elles-mêmes.
À la fin il peut y avoir une mise en pratique plus poussée.

% TODO: durée approximative


\section{Structure}


\subsection{Déplacements -- leçons du kung-fu}


\auteur{Thomas Mainguy}


\begin{enumerate}
	\item Rappels
	\begin{itemize}
		\item marche : la jambe arrière pousse ;
		\item quatre changements de garde : avant/arrière, sur les côtés (page~\pageref{dep:def:changement-garde-avant}).
	\end{itemize}
	
	\item Décomposition du changement de garde en deux étapes (page~\pageref{def:texte:garde-kung-fu}) → plus de flexibilité (peut changer de direction au milieu).
	
	\item Exercice : pratique des différents changements
	\begin{itemize}
		\item avant (technique~\ref{att:tech:changement-garde-2-temps-avant}) : 1) pivot des hanches, poids sur la jambe avant, jambe arrière commence à avancer, 2) fin du mouvement ;
		\item latéral (technique~\ref{att:tech:changement-garde-2-temps-latéral}) : 1) pivot des hanches pour amener la jambe arrière sur le côté, 2) fin du mouvement.
	\end{itemize}
	
	\item Exercice~\ref{att:ex:marche-poing-détendu} : marche avant avec coup de poing, revenir détendu en garde pour absorber l'impact.
	
	\item Ajout du haut du corps : sentir que sur un changement de garde on peut effectuer deux mouvements.
	
	\item Exercice~\ref{att:ex:changement-garde-2-temps-vide} : faire des séries de changement de garde dans le vide.
	
	\item Exercice~\ref{att:ex:changement-garde-2-temps-test} : \A et \D sont légèrement hors distance, \A effectue le premier mouvement du changement de garde.
	Si \A n'est pas bien protégé \D attaque, sinon \A termine son mouvement.
	
	\item Exercice~\ref{att:ex:changement-garde-2-temps-projection} : \A donne un coup de poing à \D qui couvre du bras en avançant.
	\D passe la jambe derrière \D tandis que la main droite pousse sur la poitrine (projection).
	
	\item Techniques~\ref{att:tech:changement-garde-2-temps-prime-haute} et \ref{att:tech:changement-garde-2-temps-latéral-prime}
\end{enumerate}


\section{Épée longue}


\subsection{Variations sur la distance}


\base{Cor Kronenburg~\cite{kronenburg:dijon:going_distance:2015}}

L'idée est d'exécuter toujours les deux mêmes mouvements – oberhau à droite puis attaque à gauche – en adaptant les détails en fonction de la distance.
Il s'agit d'une étude visant à améliorer la prise de décision.
De plus cela permet d'unifier plusieurs techniques de Liechtenauer qui ont l'air différentes prises une par une.


\begin{itemize}
	\item Notions : distances en escrime allemande : zufechten, fechten, kriegen, ringen (section~\ref{chap:épée-longue-liechtenauer}).
	\item Exercices~\ref{struct:ex:contact:frappe-épaules}, \ref{struct:ex:frappe-gauche-droite}
	\item Exercices~\ref{ex:frappe-dist:approche-croix-aleat}, \ref{ex:frappe-dist:approche-croix-aleat-garde} : \D avance et place une cible avec la garde de son épée, \A porte deux attaques, une dans chaque angle ; sans puis avec retour en garde de \A.
	\item Exercice~\ref{ex:frappe-dist:approche-double-frappe} (à gauche en changeant de pied) : \D approche, \A attaque quand il est à distance, \D recule et \A attaque à nouveau.
	\item Techniques~\ref{épée-longue:tech:dg-zufechten-abnemmen}, \ref{épée-longue:tech:dg-fechten-duplieren}, \ref{épée-longue:tech:dg-kriegen-absetzen}, \ref{épée-longue:tech:dg-armringen-einlauffen}, \ref{épée-longue:tech:dg-leibringen-einlauffen} : abnemmen, duplieren, absetzen, et einlauffen (simple et projection).
	\item Pratique : tout le monde en colonne face à \D, une personne \A s'avance et attaque, \D choisit une des cinq réactions précédentes, et \A doit réagir correctement.
\end{itemize}


\section{Messer}


\subsection{Principes internes}


\base{Martin Enzi~\cite{enzi:dijon:messer_inner:2015}}

Il s'agit de développer les concepts décrits dans la section~\ref{sec:structure:général} à l'aide des exercices~\ref{mains-nues:ex:enzi-1}, \ref{mains-nues:ex:enzi-2}, \ref{mains-nues:ex:enzi-3} et \ref{mains-nues:ex:enzi-4}, afin d'appliquer des principes au messer (section~\ref{sec:messer:lekuchner}).
L'idée est de travailler d'abord à mains nues pour voir le réflexe naturel, et ensuite ajouter l'arme.



\printbibliography[heading=bibintoc]

\end{document}
